% ══════════════════════════════════════════════════════════════
%  Section 9.3: Chain State and Storage --- How Blocks Become Persistent State
% ══════════════════════════════════════════════════════════════
\chapter{Chain: Chain State and Storage}
\label{ch:chain-state}

\begin{learnbox}
\begin{itemize}
  \item How blocks are persisted to sled, the embedded key-value database
  \item The ``tip pointer'' pattern: where the current chain head lives in-memory and on-disk
  \item Entry points for reads (status queries) and writes (appending blocks)
  \item How the chain is initialized (genesis bootstrap) and reopened from disk
  \item Iterator pattern: walking the chain backward from tip to genesis
\end{itemize}
\end{learnbox}

\begin{infobox}[\index{peer-to-peer}Distributed systems note]
In a \index{blockchain}blockchain network, each node maintains its own local copy of the \index{blockchain}blockchain. When we discuss ``chain state and storage'' in this section, we are describing how one node stores \index{block}blocks locally, not a shared database that all nodes access. Each node's \code{BlockchainFileSystem} persists \index{block}blocks to its own local \index{sled}sled database, and nodes synchronize with each other by exchanging \index{block}blocks over the network. This decentralization is fundamental to \index{blockchain}blockchain's resilience: each node operates independently with its own copy of the chain.
\end{infobox}

% ──────────────────────────────────────────────────────────────
\section{Status Read and Append-Block Write Entry Points}
\label{sec:ch09-3-entry-points}

\subsection{Entry point 1: Status read (GET blockchain info)}
\label{sec:ch09-3-status-read}

\begin{listing}[H]
\caption{Status read handler: retrieve height and tip block}
\label{lst:ch09-3-1}
\begin{minted}[fontsize=\footnotesize]{rust}
// Source: blockchain.rs --- router -> get_blockchain_info(handler)
pub async fn get_blockchain_info(
    State(node): State<Arc<NodeContext>>,
) -> Result<Json<ApiResponse<BlockchainInfoResponse>>, StatusCode> {
    let height = node
        .get_blockchain_height()
        .await
        .map_err(|_| StatusCode::INTERNAL_SERVER_ERROR)?;

    let last_block = node
        .blockchain()
        .get_last_block()
        .await?;
    let last_block_hash = last_block
        .map(|block| block.get_hash().to_string())
        .unwrap_or_else(|| "genesis".to_string());

    let mempool_size = node.get_mempool_size()?;
    // ... build response ...
    Ok(Json(ApiResponse::success(info)))
}
\end{minted}
\end{listing}

\subsection{Entry point 2: Append-block write (POST mine to address)}
\label{sec:ch09-3-append-write}

\begin{listing}[H]
\caption{Mining write handler: mine and append a block}
\label{lst:ch09-3-2}
\begin{minted}[fontsize=\footnotesize]{rust}
// Source: mining.rs --- router -> generate_to_address(handler)
let mined_block = node
    .mine_block(&transactions)
    .await
    .map_err(|_| StatusCode::INTERNAL_SERVER_ERROR)?;
\end{minted}
\end{listing}

% ──────────────────────────────────────────────────────────────
\section{Height Read: The Shared-Lock Pattern Through Storage}
\label{sec:ch09-3-height-read}

The call chain for a height read goes: Axum router \(\to\) \code{NodeContext::get\_blockchain\_height} \(\to\) \code{BlockchainService::get\_best\_height} (shared lock) \(\to\) \code{BlockchainFileSystem::get\_best\_height}.

\begin{listing}[htbp]
\caption{\code{NodeContext::get\_blockchain\_height} delegates to service}
\label{lst:ch09-3-3}
\begin{minted}[fontsize=\footnotesize]{rust}
// Source: context.rs
pub async fn get_blockchain_height(&self) -> Result<usize> {
    self.blockchain.get_best_height().await
}
\end{minted}
\end{listing}

\begin{listing}[htbp]
\caption{\code{BlockchainService::get\_best\_height} acquires shared lock}
\label{lst:ch09-3-4}
\begin{minted}[fontsize=\footnotesize]{rust}
// Source: chainstate.rs
pub async fn get_best_height(&self) -> Result<usize> {
    self.read(|blockchain: BlockchainFileSystem| async move {
        blockchain.get_best_height().await
    })
    .await
}
\end{minted}
\end{listing}

\begin{listing}[htbp]
\caption{\code{BlockchainFileSystem::get\_best\_height} hits sled}
\label{lst:ch09-3-5}
\begin{minted}[fontsize=\footnotesize]{rust}
// Source: file_system_db_chain.rs
pub async fn get_best_height(&self) -> Result<usize> {
    if self.is_empty() {
        Ok(0)
    } else {
        let block_tree = self.blockchain.db
            .open_tree(self.get_blocks_tree_path())?;
        let tip_block_bytes = block_tree
            .get(self.get_tip_hash().await?)?
            .ok_or(BtcError::BlockNotFound)?;
        let tip_block = Block::deserialize(
            tip_block_bytes.as_ref()
        )?;
        Ok(tip_block.get_height())
    }
}
\end{minted}
\end{listing}

\begin{notebox}
\code{BlockchainService::get\_best\_height(...)} is a \textbf{shared-lock read}: it uses \code{read(...)} (shared lock) and delegates to \code{BlockchainFileSystem}. The persistence engine hits sled via \code{self.blockchain.db} and uses the current tip hash to load the tip block bytes.
\end{notebox}

% ──────────────────────────────────────────────────────────────
\section{Tip-Block Read: Cache Then Lookup}
\label{sec:ch09-3-tip-read}

\begin{listing}[H]
\caption{\code{BlockchainFileSystem::get\_last\_block} and cached tip pointer}
\label{lst:ch09-3-6}
\begin{minted}[fontsize=\footnotesize]{rust}
// Source: file_system_db_chain.rs
pub async fn get_last_block(&self) -> Result<Option<Block>> {
    let tip_hash = self.get_tip_hash().await?;
    let block = self.get_block(tip_hash.as_bytes()).await?;
    Ok(block)
}

pub async fn get_tip_hash(&self) -> Result<String> {
    let tip_hash = self.blockchain.tip_hash.read().await;
    Ok(tip_hash.clone())
}
\end{minted}
\end{listing}

\begin{checkpointbox}
This is a ``tip pointer + lookup'' pattern: read the cached tip hash from the in-memory lock, then use it as a sled key to load block bytes, then deserialize.
\end{checkpointbox}

% ──────────────────────────────────────────────────────────────
\section{Where Chain State Lives: The Storage Container}
\label{sec:ch09-3-storage-container}

\begin{listing}[H]
\caption{\code{BlockchainFileSystem} owns \code{Blockchain<Db>}}
\label{lst:ch09-3-7}
\begin{minted}[fontsize=\footnotesize]{rust}
// Source: file_system_db_chain.rs
#[derive(Clone, Debug)]
pub struct BlockchainFileSystem {
    blockchain: Blockchain<Db>,
    file_system_tree_dir: String,
}

// Source: blockchain.rs --- chain container
pub struct Blockchain<T> {
    pub tip_hash: Arc<TokioRwLock<String>>,
    pub db: T,                              // sled Db handle
    pub is_empty: bool,
}
\end{minted}
\end{listing}

% ──────────────────────────────────────────────────────────────
\section{Create and Open: Bootstrap vs.\ Reopen}
\label{sec:ch09-3-create-open}

\subsection{Constructors: initialize vs.\ default}
\label{sec:ch09-3-constructors}

\begin{listing}[H]
\caption{\code{BlockchainService} constructor entry points}
\label{lst:ch09-3-8}
\begin{minted}[fontsize=\footnotesize]{rust}
// Source: chainstate.rs
impl BlockchainService {
    pub async fn initialize(
        genesis_address: &WalletAddress,
    ) -> Result<BlockchainService> {
        let blockchain =
            BlockchainFileSystem::create_blockchain(
                genesis_address
            ).await?;
        Ok(BlockchainService(Arc::new(
            TokioRwLock::new(blockchain)
        )))
    }

    pub async fn default() -> Result<BlockchainService> {
        let blockchain =
            BlockchainFileSystem::open_blockchain().await?;
        Ok(BlockchainService(Arc::new(
            TokioRwLock::new(blockchain)
        )))
    }
}
\end{minted}
\end{listing}

\begin{notebox}
\code{initialize(...)} delegates to \code{BlockchainFileSystem::create\_blockchain(...)} (genesis bootstrapping lives there). \code{default()} delegates to \code{BlockchainFileSystem::open\_blockchain(...)} (no genesis; fail if missing tip key).
\end{notebox}

\subsection{Persisted Keys That Define ``Chain Exists''}
\label{sec:ch09-3-persisted-keys}

\begin{listing}[H]
\caption{Storage constants for tip pointer and blocks}
\label{lst:ch09-3-9}
\begin{minted}[fontsize=\footnotesize]{rust}
// Source: file_system_db_chain.rs
// persisted head pointer (sled key)
const DEFAULT_TIP_BLOCK_HASH_KEY: &str = "tip_block_hash";
// sled tree: "<block_hash>" -> bytes(Block)
const DEFAULT_BLOCKS_TREE: &str = "blocks1";
// on-disk directory for the default sled DB
const DEFAULT_TREE_DIR: &str = "data1";
\end{minted}
\end{listing}

\begin{notebox}
The persisted head pointer is \textbf{one key}: \code{"tip\_block\_hash"}. The historical log is a \textbf{tree} that maps \code{<block\_hash> -> serialized Block bytes}.
\end{notebox}

\subsection{Open path: load tip from sled}
\label{sec:ch09-3-open-path}

\begin{listing}[H]
\caption{\code{BlockchainFileSystem::open\_blockchain}}
\label{lst:ch09-3-10}
\begin{minted}[fontsize=\footnotesize]{rust}
// Source: file_system_db_chain.rs
pub async fn open_blockchain()
    -> Result<BlockchainFileSystem>
{
    let db = sled::open(path)?;
    let blocks_tree = db
        .open_tree(file_system_tree_dir.clone())?;

    let tip_bytes = blocks_tree
        .get(DEFAULT_TIP_BLOCK_HASH_KEY)?
        .ok_or(BtcError::BlockchainNotFoundError(
            "No existing blockchain found.".to_string(),
        ))?;

    let tip_hash = String::from_utf8(tip_bytes.to_vec())?;

    Ok(BlockchainFileSystem {
        blockchain: Blockchain {
            tip_hash: Arc::new(TokioRwLock::new(tip_hash)),
            db,
            is_empty: false,
        },
        file_system_tree_dir,
    })
}
\end{minted}
\end{listing}

\begin{checkpointbox}
The \textbf{only ``open'' state we must recover} is the tip hash (\code{tip\_block\_hash} key). After this returns, reads like \code{get\_tip\_hash()} do \textbf{not} need to hit sled for the tip pointer (they read the in-memory lock).
\end{checkpointbox}

\subsection{Create path: bootstrap genesis if missing}
\label{sec:ch09-3-create-path}

\begin{listing}[H]
\caption{\code{BlockchainFileSystem::create\_blockchain}}
\label{lst:ch09-3-11}
\begin{minted}[fontsize=\footnotesize]{rust}
// Source: file_system_db_chain.rs
pub async fn create_blockchain(
    genesis_address: &WalletAddress,
) -> Result<Self> {
    let db = sled::open(path)?;
    let blocks_tree =
        db.open_tree(file_system_tree_dir.clone())?;

    let data = blocks_tree.get(DEFAULT_TIP_BLOCK_HASH_KEY)?;

    let tip_hash = if let Some(data) = data {
        String::from_utf8(data.to_vec())?
    } else {
        let coinbase_tx = Transaction::new_coinbase_tx(
            genesis_address
        )?;
        let block = Block::generate_genesis_block(&coinbase_tx);
        Self::update_blocks_tree(&blocks_tree, &block).await?;
        String::from(block.get_hash())
    };

    Ok(BlockchainFileSystem {
        blockchain: Blockchain {
            tip_hash: Arc::new(TokioRwLock::new(tip_hash)),
            db,
            is_empty: false,
        },
        file_system_tree_dir,
    })
}
\end{minted}
\end{listing}

\begin{notebox}
Genesis is created \textbf{exactly once} (when the tip key is absent). The persisted ``chain exists'' signal is: the \code{"tip\_block\_hash"} key is present.
\end{notebox}

% ──────────────────────────────────────────────────────────────
\section{Atomic Persistence Unit: Block Bytes + Tip Key}
\label{sec:ch09-3-atomic-unit}

\begin{listing}[H]
\caption{Atomic persistence: write block and tip together}
\label{lst:ch09-3-12}
\begin{minted}[fontsize=\footnotesize]{rust}
// Source: file_system_db_chain.rs
async fn update_blocks_tree(
    blocks_tree: &Tree,
    block: &Block
) -> Result<()> {
    let block_hash = block.get_hash();
    let block_ivec = IVec::try_from(block.clone())?;

    let transaction_result = blocks_tree.transaction(|tx_db| {
        tx_db.insert(block_hash, block_ivec.clone())?;
        tx_db.insert(DEFAULT_TIP_BLOCK_HASH_KEY, block_hash)?;
        Ok(())
    });

    transaction_result
        .map(|_| ())
        .map_err(|e| BtcError::BlockchainDBconnection(
            format!("{:?}", e)
        ))
}
\end{minted}
\end{listing}

\begin{notebox}
``Atomic'' in this context means: \textbf{block bytes plus tip pointer} are updated together inside one sled transaction. Derived state (UTXO set) is \textbf{not} part of this transaction (it is updated as a later step).
\end{notebox}

% ──────────────────────────────────────────────────────────────
\section{Mining Boundary: Append Block via Local Mining}
\label{sec:ch09-3-mining-append}

\begin{listing}[H]
\caption{Mining: verify before acquiring write lock}
\label{lst:ch09-3-13}
\begin{minted}[fontsize=\footnotesize]{rust}
// Source: chainstate.rs
pub async fn mine_block(
    &self,
    transactions: &[Transaction]
) -> Result<Block> {
    for tx in transactions {
        if !tx.verify(self).await? {
            return Err(BtcError::InvalidTransaction);
        }
    }
    let blockchain_guard = self.0.write().await;
    blockchain_guard.mine_block(transactions).await
}
\end{minted}
\end{listing}

\begin{listing}[htbp]
\caption{Persistence engine: create, persist, cache, update UTXO}
\label{lst:ch09-3-14}
\begin{minted}[fontsize=\footnotesize]{rust}
// Source: file_system_db_chain.rs
pub async fn mine_block(
    &self,
    transactions: &[Transaction]
) -> Result<Block> {
    let best_height = self.get_best_height().await?;
    let block = Block::new_block(
        self.get_tip_hash().await?,
        transactions,
        best_height + 1,
    );
    let block_hash = block.get_hash();

    let blocks_tree = self.blockchain.db
        .open_tree(self.get_blocks_tree_path())?;
    Self::update_blocks_tree(&blocks_tree, &block).await?;

    self.set_tip_hash(block_hash).await?;
    self.update_utxo_set(&block).await?;

    Ok(block)
}
\end{minted}
\end{listing}

\begin{notebox}
\code{BlockchainService::mine\_block(...)} is the \textbf{exclusive lock boundary} for local chain extension. Persistence is split into:
\begin{itemize}
  \item \textbf{Atomic persisted unit}: write \code{<block\_hash> -> bytes(Block)} and \code{"tip\_block\_hash" -> <block\_hash>}
  \item \textbf{Runtime cache update}: \code{set\_tip\_hash(...)} keeps the in-memory cached tip consistent with the persisted value
  \item \textbf{Derived state update}: UTXO updates happen after the block is persisted
\end{itemize}
\end{notebox}

% ──────────────────────────────────────────────────────────────
\section{Chain Iteration: Walking Backward from Tip to Genesis}
\label{sec:ch09-3-iteration}

\begin{listing}[H]
\caption{Iterator entry point (shared lock)}
\label{lst:ch09-3-15}
\begin{minted}[fontsize=\footnotesize]{rust}
// Source: chainstate.rs
pub async fn iterator(&self) -> Result<BlockchainIterator> {
    self.read(|blockchain| async move {
        blockchain.iterator().await
    })
    .await
}
\end{minted}
\end{listing}

\begin{listing}[htbp]
\caption{Iterator creation and stepping}
\label{lst:ch09-3-16}
\begin{minted}[fontsize=\footnotesize]{rust}
// Source: file_system_db_chain.rs
pub async fn iterator(&self) -> Result<BlockchainIterator> {
    let hash = self.get_tip_hash().await?;
    Ok(BlockchainIterator::new(
        hash,
        self.blockchain.db.clone(),
        self.get_blocks_tree_path(),
    ))
}

pub struct BlockchainIterator {
    db: Db,
    file_system_blocks_tree: String,
    current_hash: String,
}

impl BlockchainIterator {
    fn new(
        tip_hash: String,
        db: Db,
        file_system_blocks_tree: String,
    ) -> Self {
        BlockchainIterator {
            current_hash: tip_hash,
            file_system_blocks_tree,
            db,
        }
    }

}

impl Iterator for BlockchainIterator {
    type Item = Block;

    fn next(&mut self) -> Option<Self::Item> {
        let block_tree = self.db
            .open_tree(self.file_system_blocks_tree.clone())
            .ok()?;

        let data = match block_tree.get(self.current_hash.clone()) {
            Ok(Some(d)) => d,
            Ok(None) => return None,
            Err(_) => return None,
        };

        let block = Block::deserialize(
            data.to_vec().as_slice()
        ).ok()?;
        self.current_hash = block.get_pre_block_hash().clone();
        Some(block)
    }
}
\end{minted}
\end{listing}

\begin{listing}[htbp]
\caption{Higher-level API: blocks by height range}
\label{lst:ch09-3-17}
\begin{minted}[fontsize=\footnotesize]{rust}
// Source: chainstate.rs
pub async fn get_blocks_by_height(
    &self,
    initial_height: usize,
    height: usize,
) -> Result<Vec<Block>> {
    let mut blocks = Vec::new();
    let iterator = self.iterator().await?;
    for block in iterator {
        let h = block.get_height();
        if h > height || h < initial_height {
            break;
        }
        blocks.push(block);
    }
    Ok(blocks)
}
\end{minted}
\end{listing}

\begin{notebox}
The iterator is exposed at the \textbf{API boundary} (\code{BlockchainService::iterator}) and implemented in the persistence engine. Iteration starts at the current tip hash and repeatedly follows \code{pre\_block\_hash} ``backward'' through history.
\end{notebox}

% ──────────────────────────────────────────────────────────────
\section{Network Insertion: Add Block from Peer}
\label{sec:ch09-3-network-insert}

\begin{listing}[H]
\caption{Network block accepted into chain}
\label{lst:ch09-3-18}
\begin{minted}[fontsize=\footnotesize]{rust}
// Source: chainstate.rs
pub async fn add_block(&self, block: &Block) -> Result<()> {
    let mut blockchain_guard = self.0.write().await;
    blockchain_guard.add_block(block).await
}
\end{minted}
\end{listing}

\begin{listing}[htbp]
\caption{Persistence layer: idempotence check, serialize, transaction, height comparison}
\label{lst:ch09-3-19}
\begin{minted}[fontsize=\footnotesize]{rust}
// Source: file_system_db_chain.rs
pub async fn add_block(&mut self, new_block: &Block) -> Result<()> {
    let block_tree = self.blockchain.db
        .open_tree(self.get_blocks_tree_path())?;

    if block_tree.get(new_block.get_hash())?.is_some() {
        return Ok(());
    }

    let block_bytes = new_block.serialize()?;
    let tip_hash = self.get_tip_hash().await?;

    let transaction_result = block_tree.transaction(|tx| {
        tx.insert(new_block.get_hash(), block_bytes.clone())?;
        if new_block.get_height() > /* tip height */ {
            tx.insert(
                DEFAULT_TIP_BLOCK_HASH_KEY,
                new_block.get_hash()
            )?;
        }
        Ok(())
    });

    transaction_result
        .map_err(|e| BtcError::BlockchainDBconnection(
            format!("{:?}", e)
        ))?;

    Ok(())
}
\end{minted}
\end{listing}

\begin{notebox}
The network-received block is persisted \textbf{before} the full ``validate \(\to\) connect'' gate is enforced (that hard gate is still incomplete in this learning implementation). Tip updates use a simplified height-based fork-choice rule.
\end{notebox}

% ──────────────────────────────────────────────────────────────
\section{Summary}
\label{sec:ch09-3-summary}

\begin{chaptersummary}
\item \textbf{Two entry points}: status reads (shared lock), append-block writes (exclusive lock)
\item \textbf{Tip pointer pattern}: in-memory cache of the canonical head hash, backed by the sled key \code{"tip\_block\_hash"}
\item \textbf{Genesis bootstrapping}: created exactly once when the tip key is absent
\item \textbf{Atomic persistence unit}: block bytes and tip pointer written together in one sled transaction
\item \textbf{Chain iteration}: walk backward from tip by following \code{pre\_block\_hash} links
\item \textbf{Network insertion}: accept candidate blocks, update tip if candidate wins by height
\end{chaptersummary}
